\chapter{Methodology}
\label{chap:methodology}

This chapter specifies the complete chain from literature evidence to held-out
predictions. The unit recorded in the database is an experimental condition,
but the unit of scientific transfer is a publication: rows from one paper share
apparatus, materials, measurement practice, and often a common experimental
campaign. The research therefore treats row-level interpolation and
paper-level extrapolation as distinct estimands, confines every data-adaptive
operation to training partitions, and preserves a machine-readable audit trail
from each reported score to its source rows, split assignment, candidate
configuration, and fitted model. These principles were fixed before comparing
candidate performance.

\section{Assembly of the microgravity flammability database}
\label{sec:database-assembly}

\subsection{Literature screening and inclusion}

The database was assembled before either machine-learning pipeline was designed.
Approximately 130 experimental papers and technical reports published from 1966
onward were screened for solid-fuel combustion under microgravity,
partial-gravity, or spacecraft-relevant atmospheric and flow conditions. Only
about 87 publications exposed sufficiently resolved observations to map ignition
or flame-spread rate (FSR) into individual database rows. A paper is useful for
modeling only when the experiment, specimen, environment, and measured outcome
could be connected at row level; a qualitative conclusion without recoverable
conditions could support the literature review but not supervised learning.

Each row retains a full human-readable  Article (MLA), Authors, and DOI when available.
That provenance is essential because observations from one publication are
clustered and because ambiguous values must be checked against the source rather
than treated as anonymous numbers.

\begin{figure}[htbp]
  \centering
  \includegraphics[width=0.92\textwidth]{database_overview.png}
  \caption{Overview of the assembled microgravity flammability table before
  task-specific cleaning. The corpus contains 5,057 rows from 87 papers
  (median 32 rows per paper). Ignition is the majority outcome (76.6\%).
  Numeric FSR is available for 2,590 raw rows. Operating envelopes span
  oxygen mole fraction \(0.002\)--\(1\), pressure from a few kilopascals to
  several megapascals, gravity from microgravity to several \(g\), and a
  wide fuel-property range. Modeling later uses SI flame-spread units of
  \si{\metre\per\second}; the overview card reports the exploratory FSR
  range in millimetres per second.}
  \label{fig:method-overview}
\end{figure}

The overview is not a balanced design. A median of 32 rows per
paper coexists with campaigns that contribute several hundred nearby
conditions from a single rig. That source-size structure is why a random
row split is optimistic: conditions from a large paper can appear on both
sides of an interpolation split, whereas holding out that paper removes a
large, internally correlated block. The same card also shows why two
supervised tasks, not five, are feasible: ignition labels are essentially
complete, numeric FSR covers about half the table, and the remaining
flammability outputs are too sparse to train.

\subsection{Extraction strategy}

The first attempted extraction workflow was a conventional optical character
recognition (OCR) pipeline. It was not viable for this corpus. The sources range
from old scanned papers to long NASA memoranda; notation, units, table
structures, and page layouts differ; scans and plotted markers can be blurred;
and many results appear only in graphs rather than machine-readable text.
Generic OCR also confuses subscripts and symbols, separates values from their
headers, and cannot reliably infer which experimental condition belongs to a
plotted point. Applying it at scale would have created a large table with an
unacceptably hidden transcription error rate.

The adopted workflow was therefore AI-assisted and human-verified:
\begin{enumerate}
  \item the source was read to identify the experiment, sample, controlled
  conditions, reported outputs, and units;
  \item an AI extraction pass was used to recover as much structured information
  as possible for each experimental row;
  \item every extracted row was then checked against the paper by the database
  curator, and missing or incorrectly associated values were completed or
  corrected manually;
  \item units and categorical names were standardized while the original
  citation and explanatory notes were retained for traceability.
\end{enumerate}

AI was used as an extraction accelerator, not as an authority. This distinction
is important for scientific data: a fluent transcription can still attach a
correct number to the wrong series, silently flatten a logarithmic axis, or
invent precision that the graph does not contain. Manual review remained the
acceptance step. Values read from plots were stored only to the precision
supported by marker size and axis resolution; visually inferred trailing digits
were not treated as measured information. Ambiguous series assignments or units
were resolved from captions and surrounding text, or left missing when the
source did not support a defensible interpretation.

For dense plots, the points were digitized with WebPlotDigitizer
\cite{webplotdigitizer}. Axis calibration was set from the published ticks,
including explicit log-axis selection where required; each color or marker
series was traced separately; overlapping points were inspected at high zoom;
and the exported coordinates were checked against axis limits and values quoted
in the text. This procedure was preferred to AI-internal OCR because exact point
coordinates, logarithmic scaling, and overlapping markers were otherwise
unreliable.

\subsection{Row structure, outputs, and quality control}

The final schema groups fields into citation, specimen, flow, apparatus,
outputs, and notes. Inputs cover geometry, dimensions, fuel and core
thermophysical properties, oxygen concentration, diluent, gas properties,
pressure, flow, dimensionless transport groups, gravity, facility, and ignition
conditions. Flammability is documented through the five outputs introduced in
\cref{tab:flammability-variables}: ignition, flame length, FSR, heat-release
rate, and smoke/aerosols. Not every paper reports every output.

The raw table uses a category banner as its first row and field names as its
second. Values were standardized to SI units wherever the source supported a
defensible conversion: length in metres, pressure in pascals, velocity in
\si{\metre\per\second}, power in watts, time in seconds, density in
\si{\kilogram\per\metre\cubed}, conductivity in
\si{\watt\per\metre\per\kelvin}, heat capacity in
\si{\joule\per\kilogram\per\kelvin}, and FSR in
\si{\metre\per\second}. Values whose physical dimension did not match a target
field were not coerced merely to fill a cell.

{\small
\setlength{\tabcolsep}{3.5pt}
\begin{longtable}{@{}p{0.12\textwidth}p{0.34\textwidth}p{0.18\textwidth}r p{0.20\textwidth}@{}}
\caption{Complete field dictionary of the database. Completeness is the
fraction of the 5,057 assembled rows with a non-empty cell. Modeling roles
state whether the field is a predictor, a target, a grouping key, or is
excluded from supervised inputs.}
\label{tab:database-schema}\\
\toprule
Group & Excel field & Units & Compl.\ (\%)& Modeling role \\
\midrule
\endfirsthead
\caption[]{Complete field dictionary of the Excel database (continued).}\\
\toprule
Group & Excel field & Units & Compl.\ (\%)& Modeling role \\
\midrule
\endhead
\bottomrule
\endlastfoot
Citation & Article (MLA) & --- & 100.0 & grouping / excluded \\
Citation & Authors & --- & 100.0 & grouping / excluded \\
Citation & DOI & --- & 93.2 & grouping / excluded \\
Sample & Geometry of Sample & categorical & 100.0 & predictor \\
Sample & Dimensions of sample & m (geometry-specific) & 100.0 & parsed into length scales \\
Sample & \texttt{half\_thickness} & m & 100.0 & predictor \\
Sample & \texttt{characteristic\_length\_m} & m & 100.0 & predictor \\
Sample & Material of sample & free text & 100.0 & excluded \\
Sample & \texttt{fuel\_density\_kg\_m3} & \si{\kilogram\per\metre\cubed} & 100.0 & predictor \\
Sample & \texttt{fuel\_k\_W\_mK} & \si{\watt\per\metre\per\kelvin} & 100.0 & predictor \\
Sample & \texttt{fuel\_cp\_J\_kgK} & \si{\joule\per\kilogram\per\kelvin} & 100.0 & predictor \\
Sample & \texttt{fuel\_pyrolysis\_T\_K} & K & 100.0 & predictor \\
Sample & \texttt{fuel\_alpha\_m2\_s} & \si{\metre\squared\per\second} & 100.0 & derived predictor \\
Sample & \texttt{fuel\_volumetric\_heat\_capacity\_J\_m3K} & \si{\joule\per\metre\cubed\per\kelvin} & 100.0 & derived predictor \\
Sample & \texttt{core\_density\_kg\_m3} & \si{\kilogram\per\metre\cubed} & 100.0 & predictor \\
Sample & \texttt{core\_k\_W\_mK} & \si{\watt\per\metre\per\kelvin} & 100.0 & predictor \\
Sample & \texttt{core\_cp\_J\_kgK} & \si{\joule\per\kilogram\per\kelvin} & 100.0 & predictor \\
Sample & \texttt{core\_volumetric\_heat\_capacity\_J\_m3K} & \si{\joule\per\metre\cubed\per\kelvin} & 100.0 & derived predictor \\
Flow & Oxygen Concentration & mole fraction & 100.0 & predictor \\
Flow & diluent & categorical & 100.0 & predictor \\
Flow & \texttt{gas\_M} & \si{\kilogram\per\mole} & 100.0 & predictor \\
Flow & \texttt{gas\_cp\_mass} & \si{\joule\per\kilogram\per\kelvin} & 100.0 & predictor \\
Flow & \texttt{gas\_k} & \si{\watt\per\metre\per\kelvin} & 100.0 & predictor \\
Flow & \texttt{gas\_density\_kg\_m3} & \si{\kilogram\per\metre\cubed} & 98.9 & predictor \\
Flow & \texttt{gas\_alpha\_m2\_s} & \si{\metre\squared\per\second} & 98.9 & derived predictor \\
Flow & \texttt{gas\_nu\_m2\_s} & \si{\metre\squared\per\second} & 98.9 & derived predictor \\
Flow & Reynolds & --- & 98.9 & derived predictor \\
Flow & Peclet & --- & 98.9 & derived predictor \\
Flow & Prandtl & --- & 98.9 & derived predictor \\
Flow & \texttt{thermal\_diffusion\_time} & s & 100.0 & derived predictor \\
Flow & Pressure (Pa) & Pa & 98.9 & predictor \\
Flow & Flow Velocity (m/s) & \si{\metre\per\second} & 79.5 & predictor \\
Flow & Rig Name & free text & 76.6 & excluded apparatus \\
Flow & Internal Geometry & free text & 86.9 & excluded apparatus \\
Flow & Internal Dimensions & m & 90.3 & excluded apparatus \\
Flow & Gravity & \(g/g_{\mathrm{earth}}\) & 98.9 & predictor \\
Extra & Facility & categorical & 98.9 & excluded apparatus \\
Extra & Ignition Method & categorical & 98.6 & predictor \\
Extra & Ignition Power (W) & W & 18.8 & predictor (sparse) \\
Extra & Ignition Time (s) & s & 45.6 & predictor (sparse) \\
Outputs & Ignition & Yes/No & 100.0 & classification target \\
Outputs & Flame Length & m & 4.3 & post-outcome; excluded \\
Outputs & FSR (m/s) & \si{\metre\per\second} & 51.5 & regression target \\
Outputs & HRR & W & 1.1 & post-outcome; excluded \\
Outputs & Smoke/Aerosols & --- & 1.8 & post-outcome; excluded \\
Info & Info & free text & 25.5 & notes; excluded \\
\end{longtable}
}

Ignition labels, fuel thermophysical properties, oxygen, and diluent are
essentially complete. Geometry, gravity, pressure, and ignition method remain
above 98\%. Flow velocity drops to 79.5\%. Numeric FSR is present in only
51.5\% of the raw rows, which is why the regression cohort is smaller than
the ignition cohort. Flame length (4.3\%), heat-release rate (1.1\%), and
smoke/aerosols (1.8\%) are too sparse to support a third supervised task.
Ignition power (18.8\%) and ignition time (45.6\%) are retained as predictors
where reported, with fold-local imputation for the remainder. Apparatus
names, facility labels, free-text material names, and notes are excluded from
model inputs because they fingerprint a laboratory without encoding
physics. The schema therefore justifies both the two-task scope and the
need for fold-local imputation: missingness is a property of what
laboratories reported, not a random mask.

Quality control is applied in the code : empty rows and
columns removed, target labels normalized, exact modeling duplicates removed,
and unparseable values recorded as missing rather than guessed. Physical-domain
checks rejected impossible concentrations, non-positive pressures or FSR; conversions were performed before duplicate detection so that
equivalent observations expressed in different units could not survive as
independent rows. Canonical paper identity is a citation. This produces identifiers of the form \texttt{citation::<normalized-citation>}.
Dataset shape,
column roles, target prevalence, paper counts, and a content fingerprint are
written to validation records and checked again when frozen splits are loaded.

The two tasks use different valid subsets. After feature/target
standardization and deduplication, ignition classification contains 4,487 rows
from 87 canonical papers (3,383 ignition and 1,104 no ignition). FSR regression
contains 2,531 rows with a numeric FSR from 66 papers. The FSR papers contribute
between 1 and 193 rows, with median 22.5. Dataset fingerprints are persisted with
the splits and evaluations so that a result cannot silently refer to a changed
database.

\subsection{Sampling Bias \& Class Bias}
\label{sec:method-imbalance}

The assembled literature table is not balanced. Two
imbalances dominate every subsequent modeling choice.

The first is source-size imbalance. Publications contribute very different
numbers of points from a single rig or campaign.
\Cref{fig:method-overview} already shows that the median paper is modest
while a handful of campaigns dominate the row count; the Pareto curve in
\cref{fig:method-imbalance}(a)
makes the cumulative effect explicit. After sorting papers from largest to
smallest contribution, 36 of 87 papers account for more than 80\% of all rows.
The remaining long tail contributes the last 20\%. A loss that treats every row
as an independent draw therefore over-represents densely sampled campaigns and
under-represents sparse but physically distinct sources. Paper weighting,
paper-grouped evaluation, and paper-cluster bootstrap intervals are the
methodological response to this right-skewed structure.\footnote{A right-skewed distribution is a distribution with a long tail toward larger values, indicating that most observations are concentrated at lower values while a smaller number of observations take substantially larger values.}

The second is class imbalance on the ignition label. In the raw assembled
table of 5,057 rows, 3,875 observations (76.6\%) are ignition and 1,182
(23.4\%) are no ignition (\cref{fig:method-imbalance}(b)). After cleaning the modeling dataset remains majority-positive: 3,383
ignition versus 1,104 no ignition. A classifier that always predicts ignition
would therefore already achieve about three-quarters accuracy and a
precision--recall baseline equal to the positive prevalence. Accuracy is not a
defensible selection metric. The ignition pipeline instead compares models with
ROC--AUC, PR--AUC, and Brier score, and it reports thresholded sensitivity,
specificity, precision, F1, balanced accuracy, and Matthews correlation. Class
weights and paper weights are tested as candidate recipes rather than assumed
to correct the imbalance automatically.

\begin{figure}[htbp]
  \centering
  \begin{subfigure}[t]{0.61\textwidth}
    \includegraphics[width=\linewidth]{database_pareto.png}
    \caption{Cumulative data share after sorting papers by contribution.
    Thirty-six papers supply more than 80\% of the rows.}
  \end{subfigure}\hfill
  \begin{subfigure}[t]{0.36\textwidth}
    \includegraphics[width=\linewidth]{database_ignition_balance.png}
    \caption{Raw ignition outcome composition: 76.6\% ignition versus
    23.4\% no ignition.}
  \end{subfigure}
  \caption{Two sampling imbalances in the assembled database. Panel (a) shows
  that a minority of publications dominate the row count; panel (b) shows that
  ignition is the majority class. Both motivate grouped evaluation and
  prevalence-aware classification metrics.}
  \label{fig:method-imbalance}
\end{figure}

\section{Methods common to both machine-learning pipelines}
\label{sec:common-pipeline}

The two pipelines share a scientific design even though one classifies
ignition and the other regresses flame-spread rate. Both estimate a mapping
from measured physical inputs \(\mathbf{x}\) to an observed target \(y\): the
classifier estimates \(\widehat p=P(y_{\mathrm{ign}}=1\mid\mathbf{x})\); the
regressor estimates \(\widehat v_{\mathrm{FSR}}\). Both freeze splits before
any model is compared, nest hyperparameter search inside outer training
partitions, and report interpolation and extrapolation as different questions.

The repeated splits estimate performance for the empirical mixture of
conditions and papers represented by the database. They do not constitute
an external validation set and paper-grouped performance should
not be interpreted as guaranteed transfer beyond the measured feature envelope.
The subsections below explain why the model families were chosen, how nested
evaluation is enforced, and which preprocessing and weighting choices are
common to both tasks.

\subsection{Benchmarking motivation and model families}

The closest prior machine-learning study of flame spread is Jose Rivera et al. \cite{riveraWireML}. They compiled approximately 1,200 electrical-wire observations and
compared linear regression, elastic net, stochastic gradient descent, Bayesian
ridge, kernel ridge, decision trees, \(k\)-nearest neighbours, gradient
boosting, support-vector regression, and a multilayer perceptron under random
10-fold cross-validation. That paper is the explicit
benchmark for family choice in this report. Nonlinear models were competitive
there; linear and shrinkage families lagged. The present work therefore
retains the families that were scientifically informative in that study,
upgrades gradient boosting to XGBoost, and adds random forest as a bagging
control for FSR. Linear, elastic-net, SGD, Bayesian-ridge, and kernel-ridge
regressors are not re-run: Jose already showed that they are a weak hypothesis
for heterogeneous flame-spread tables, and repeating them would inflate the
portfolio without a new physical question.

XGBoost is the primary tabular learner because Shwartz-Ziv and Armon found
gradient-boosted trees to be a strong, readily tuned baseline on heterogeneous
tables, including settings where deep models did not dominate
\cite{shwartzZiv2022tabular,chenGuestrin2016}. That prior is a starting point,
not a result. Decision trees remain as a transparent nonlinear baseline from
Jose. \(k\)-NN tests whether nearby experimental conditions are locally
predictive. SVM/SVR tests linear versus radial kernels in the engineered
physics space. MLPs test whether a smooth universal approximator helps once
the same nested, leakage-controlled protocol is imposed. Random forest is the
one family Rivera et al. did not include; it isolates bagging from boosting on the FSR
task.

\subsection{Features, leakage exclusions, and preprocessing}

The two pipelines use measurements available before the outcome is known. The
ignition pipeline predicts whether ignition occurs; the FSR pipeline predicts
the measured flame-spread rate. Each task therefore uses a different valid
subset of the database and a task-specific feature manifest.

Both pipelines use physical descriptors of the specimen, gas, operating
environment, and ignition conditions. These include:

\begin{itemize}
  \item sample geometry and characteristic dimensions;
  \item fuel and core density, conductivity, heat capacity, pyrolysis
  temperature, thermal diffusivity, and volumetric heat capacity;
  \item oxygen concentration, diluent, gas properties, and dimensionless
  transport quantities;
  \item pressure, flow velocity, gravity, ignition method, ignition power, and
  ignition duration where available;
  \item derived physical quantities, including Reynolds number, Peclet number,
  Prandtl number, and thermal diffusion time.
\end{itemize}

The ignition pipeline uses a fixed manifest of 32 numeric and five categorical
features. The FSR pipeline detects the available numeric and categorical fields
from the same cleaned schema and excludes fields that are unavailable or
inappropriate for FSR prediction.

Several fields are deliberately excluded from both pipelines:

\begin{itemize}
  \item Paper identifiers, citation fields, DOI, author names, facility names,
  rig names, internal apparatus descriptions, free-text material labels, and
  notes are excluded because they can identify an experimental source without
  representing a transferable physical mechanism.

  \item Flame length, heat-release rate, smoke/aerosols, extinction indicators,
  and duplicate target columns are excluded because they are outcomes,
  post-outcome measurements, or direct alternatives to the target being
  predicted.

  \item The ignition label is the classification target and is not used as an
  FSR predictor. Similarly, numeric FSR is the regression target and is never
  used as an ignition predictor.
\end{itemize}

All cleaning and preprocessing operations are fitted using training data only.
The outer test partition is not used to calculate a median, scale a feature,
define a category vocabulary, estimate a sample weight, choose a threshold, or
select a hyperparameter.

The preprocessing sequence within each training partition is:

\begin{enumerate}
  \item Numeric values are parsed, converted to consistent SI units where
  appropriate, and checked for physically invalid values. Invalid or
  unparseable values are recorded as missing rather than guessed.

  \item Exact duplicate modelling rows are removed after standardisation, so
  equivalent measurements expressed in different units are not counted as
  independent observations.

  \item For model families that require complete numeric inputs, missing numeric
  values are replaced with the median calculated from the current training
  partition. XGBoost retains native numeric missing values.

  \item Missing categorical values are assigned an explicit \texttt{Unknown}
  category. Categorical variables are one-hot encoded using categories observed
  in the training partition. A category first encountered in the test partition
  does not create a new fitted column.

  \item Distance-based, margin-based, and neural-network models standardise
  numeric variables using statistics calculated from the current training
  partition. Tree-based models do not require this scaling.

  \item Paper weights, class weights, bootstrap-augmentation draws, and
  threshold-selection rules are calculated only after the outer training
  partition has been defined.
\end{enumerate}

This procedure prevents information from the outer test partition from
influencing the fitted representation. It also ensures that any performance
difference between candidates is caused by their modelling choices rather than
by a preprocessing step that has seen the test data.
\subsection{Nested evaluation}

Nested evaluation separates model fitting, model selection, and final
performance measurement. This separation is necessary to obtain a fair
estimate of how well a model is expected to perform on new data.

\begin{description}
  \item[Training data] are used to fit the model parameters. For example,
  this includes learning tree splits, regression coefficients, or neural
  network weights.

  \item[Validation data] are used to compare model choices and
  hyperparameter settings. Examples include tree depth, regularization
  strength, and the number of neighbours.

  \item[Test data] are kept completely separate until the final evaluation.
  They are not used to train a model or choose its settings.
\end{description}

The nested evaluation procedure used in this research follows these steps:

\begin{enumerate}
  \item The data are divided into an outer training partition and an outer
  test partition.

  \item Only the outer training partition is used during model development.
  The outer test partition remains untouched.

  \item Within the outer training partition, inner cross-validation compares
  candidate models and their hyperparameter settings.

  \item The candidate and setting with the best average inner-validation score
  are selected.

  \item The selected configuration is retrained using all observations in the
  outer training partition.

  \item The retrained model is evaluated once on the untouched outer test
  partition. This score is used as the performance estimate for that outer
  split.
\end{enumerate}

All preprocessing steps, including transformations and feature selection when
applicable, are performed separately within the training portion of each
split. This prevents information from the validation or test data from
influencing model selection.

The same outer splits are used for all candidate models so that their results
can be compared under identical data partitions. Repeating the outer split
also measures how sensitive the results are to the particular division of the
data.

\subsection{Interpolation versus extrapolation}

Both pipelines answer two different questions :
\begin{itemize}
  \item \textbf{Interpolation} holds out rows while allowing other rows from
  the same paper to remain in training. It asks whether the model predicts new
  conditions inside campaigns it partly knows.
  \item \textbf{Extrapolation} holds out complete papers. It asks
  whether the procedure transfers to a new publication, rig, or campaign and is
  the primary generalization test.
\end{itemize}

The distinction prevents paper-specific sampling and apparatus signatures from
being mistaken for broad physical transfer. Rivera et al.\ themselves describe
random cross-validation as robustness within a compiled table rather than
strict held-source transfer \cite{riveraWireML}. The present design makes that
caveat operational: all candidates within a task share the same frozen outer
assignments, and grouped protocols assert zero paper overlap.

\subsection{Paper weighting, uncertainty, and refitting}

During training, a model makes a prediction for each row and is penalized when
that prediction differs from the observed value. This penalty is called the
\emph{loss}. A smaller loss means a better prediction: a loss of zero would
mean a perfect prediction for that row.

The model-training objective combines these row-level losses into one value.
This value, called the \emph{empirical risk}, is simply the average loss over
the rows in the training partition. Training aims to find a model with low
empirical risk.

In this dataset, as some papers contribute many more rows than others, without weighting, a paper with many rows contributes more losses to the
average and can therefore have more influence on the fitted model. Paper
weighting changes how much each paper contributes to the training objective.
It is tested as a modeling choice rather than assumed to improve performance.

For a model \(f\), row \(i\), observed outcome \(y_i\), and prediction
\(f(\mathbf{x}_i)\), the loss
\(\ell\!\left(y_i,f(\mathbf{x}_i)\right)\) measures the prediction error for
that row. Let \(p(i)\) denote the paper associated with row \(i\), let \(n_p\)
be the number of training rows from paper \(p\), and let \(w_{p(i)}\) be the
weight assigned to that paper. The weighted empirical risk is

\[
  \widehat R_w(f)=\frac{\sum_{i\in\mathcal T}w_{p(i)}
  \ell\!\left(y_i,f(\mathbf{x}_i)\right)}
  {\sum_{i\in\mathcal T}w_{p(i)}}.
\]

In this equation, each row-level loss is multiplied by the weight of its
source paper before the losses are averaged. A larger paper weight gives that
paper more influence during training. The denominator divides by the total
weight, so the result remains an average loss that can be compared across
training folds.

Both pipelines compare three weighting strategies
(\cref{fig:method-paper-weights}):

\begin{description}
  \item[Unweighted training: \(w_p=1\).]
  Every row has the same weight. A paper with more rows contributes more total
  loss and therefore has more influence on the fitted model.

  \item[Square-root weighting: \(w_p=n_p^{-1/2}\).]
  Large papers still have more influence than small papers, but their
  advantage is reduced.

  \item[Inverse weighting: \(w_p=n_p^{-1}\).]
  Each paper has approximately the same total influence, regardless of its
  number of rows.
\end{description}

Weights and paper counts are calculated using only the current training
partition. For model families that do not directly accept sample weights, the
training data are resampled using a deterministic probability-proportional
procedure. This ensures that the selected weighting strategy is applied
consistently to every candidate model.

\begin{figure}[htbp]
  \centering
  \includegraphics[width=0.98\textwidth]{method_paper_weights.pdf}
  \caption{Paper-weighting strategies shared by ignition and FSR candidates.
  \(n_p\) is the number of training rows from paper \(p\). The preferred
  strategy is selected empirically from nested scores, not assumed a priori.}
  \label{fig:method-paper-weights}
\end{figure}
\paragraph{Uncertainty of reported performance.}

Outer-fold means and standard deviations show how much model performance
changes across different data partitions. The \(95\%\) percentile intervals are
calculated with 2,000 bootstrap replicates. The ignition-classification and FSR
regression pipelines use different bootstrap procedures because they store and
evaluate different types of outputs.
\newpage
\begin{description}
  \item[Ignition classification :]
  For each candidate and evaluation protocol, every outer fold produces one
  held-out ROC-AUC, PR-AUC, and Brier score. The bootstrap resamples these
  outer-fold metric values with replacement and calculates their mean for each
  replicate. The \(95\%\) interval is the range from the 2.5th to the 97.5th
  percentile of the 2,000 bootstrap means.

  This classifier procedure is the same for interpolation and extrapolation:
  it measures variation among repeated outer-fold scores. It does not resample
  individual observations or complete paper clusters. Threshold-dependent
  metrics, such as MCC, F1 score, balanced accuracy, sensitivity, specificity,
  and precision, are calculated with the thresholds selected during nested
  evaluation, but bootstrap intervals are not reported for those metrics.

  \item[FSR regression :]
  The FSR pipeline stores the measured and predicted flame spread rate for
  every outer-test observation. When the same observation is evaluated in more
  than one repeated outer split, its predictions are averaged first. This
  produces one paired measured and predicted FSR value for each observation.

  For interpolation, the bootstrap resamples these unique observations with
  replacement. RMSE, MAE, and \(R^2\) are recalculated from the resampled
  measured--predicted pairs for each replicate.

  For extrapolation, the bootstrap resamples complete canonical-paper clusters
  with replacement. When a paper is selected, all observations from that paper
  are included in the replicate. If the same paper is selected more than once,
  all of its observations are repeated accordingly. RMSE, MAE, and \(R^2\)
  are then recalculated from the resulting paper-cluster sample.

  The \(95\%\) FSR interval is defined by the 2.5th and 97.5th percentiles of
  the 2,000 bootstrap values for each metric.
\end{description}

The classifier intervals quantify variation among repeated outer-fold scores.
The FSR intervals quantify resampling uncertainty in the stored out-of-sample
prediction pairs. Neither procedure includes uncertainty from literature
selection, data transcription, or experimental measurement.

\paragraph{Final refitting and champion selection.}

After model selection is complete, the chosen configuration is refitted using
all labeled rows. This full-data model is the final inference artifact and can
be used for prediction.

Its training performance is not used as evidence that the model generalizes to
new data. Generalization claims are based only on the nested evaluation
results obtained before refitting.

A model can be called a champion only when it appears in the saved selection
evidence. In particular:

\begin{itemize}
  \item The merged FSR comparison includes 25 of the 34 configured
  candidates.

  \item The formal ignition selection record does not include the available
  MLP shards.
\end{itemize}

Therefore, a configured candidate cannot become the reported champion unless
it was fully evaluated and included in the persisted comparison record.
\subsection{Reproducibility and audit trail}
\label{sec:method-reproducibility}

Reproducibility is enforced through persisted artifacts rather than relying only
on narrative descriptions of the analysis. Each cleaned task-specific dataset is
fingerprinted before split generation. Outer and inner split assignments are
created from declared seeds, stored with canonical paper identifiers, and reused
by every candidate.

For every candidate and outer split, the pipeline records:

\begin{itemize}
  \item the candidate identity, feature manifest, paper-weighting policy, and
  augmentation condition;
  \item the selected hyperparameters and inner-validation history;
  \item fold-level metrics and row-level outer-test predictions;
  \item paper-level metrics where applicable;
  \item integrity checks, including expected-fold completion, finite metrics,
  valid classifier probabilities, and zero paper overlap for grouped splits.
\end{itemize}

Candidate selection reads these persisted evaluation artifacts without fitting a
new model. Candidates with incomplete folds, invalid predictions, or failed
integrity checks are excluded from the corresponding comparison. Only after
selection is complete is the chosen configuration refitted on all labelled rows
to create an inference artifact.

The evaluation artifacts and refitted artifacts have different roles.
Outer-fold predictions are the comparison evidence because they are generated
on untouched test partitions. The final refitted model is trained on all
labelled rows and is used for inference, but its training fit is not used as
evidence of generalisation.

\subsection{Explainability and physical interpretation}
\label{sec:method-explainability}

Explainability is performed after model selection. It does not influence model
ranking, candidate selection, hyperparameter optimisation, or threshold
selection. Its purpose is to examine how a fitted model uses the available
features and to identify hypotheses for future experiments.

Three complementary methods are used:

\begin{description}
  \item[Tree feature importance.]
  Tree importance measures how often and how effectively a feature is used to
  split observations within a fitted tree ensemble. It provides a fast
  model-specific ranking, but it can favour variables with many possible split
  points or variables correlated with source-specific structure.

  \item[Permutation importance.]
  Permutation importance measures the loss of held-out performance after one
  original feature is shuffled while all other features remain unchanged. For
  ignition, the loss is measured using held-out ROC--AUC. For FSR, it is
  measured using the increase in held-out RMSE. A large performance loss means
  that the fitted model relied on that feature for the evaluated data.

  \item[SHAP values.]
  SHAP decomposes a fitted prediction into a baseline value and additive
  feature contributions:
  \[
  f(\mathbf{x}_i)
  =
  \phi_0+\sum_{j=1}^{p}\phi_{ij},
  \]
  where \(f(\mathbf{x}_i)\) is the prediction for observation \(i\),
  \(\phi_0\) is the model baseline, and \(\phi_{ij}\) is the contribution of
  feature \(j\) to that prediction.
\end{description}

The global mean absolute SHAP importance of feature \(j\) is

\[
I^{\mathrm{SHAP}}_j
=
\frac{1}{n}\sum_{i=1}^{n}|\phi_{ij}|.
\]

A large value of \(I^{\mathrm{SHAP}}_j\) means that the feature changes model
predictions strongly on average. It does not indicate whether the feature
increases or decreases the prediction.

SHAP beeswarm plots retain this missing information. Each point represents one
observation. Its horizontal position is the SHAP value: positive values push a
prediction upward, while negative values push it downward. For ignition,
positive values push the predicted probability toward ignition and negative
values push it toward no ignition. For FSR, positive values push the predicted
flame-spread rate upward and negative values push it downward. Point colour
shows the original feature value: red indicates a high feature value and blue
indicates a low feature value.

These methods explain the fitted predictive function; they do not establish
causal combustion mechanisms. Correlated raw and derived variables can
substitute for one another, missingness can encode source structure, and a
physically plausible feature ranking does not prove that changing that feature
would cause the predicted outcome to change. Explainability results are
therefore interpreted as model-specific hypotheses and experimental
priorities, not as universal causal rankings.
\section{Ignition-classification pipeline}
\label{sec:classification-method}
\subsection{Task, probabilities, and metrics}

Ignition is a binary target: each observation is labelled ignition
(\(y=1\)) or no ignition (\(y=0\)). The classifier does not directly return a
physical yes/no switch. Instead, it returns an estimated ignition probability
\(\widehat p\in[0,1]\).

A decision threshold \(\tau\) converts this probability into a binary
prediction:
\[
\widehat y(\tau)=
\begin{cases}
1, & \widehat p\geq\tau,\\
0, & \widehat p<\tau.
\end{cases}
\]

A lower threshold identifies more possible ignitions, but usually produces more
false alarms. A higher threshold reduces false alarms, but increases the risk
of missing true ignitions. Each threshold therefore produces a different
combination of true positives (TP), false positives (FP), true negatives (TN),
and false negatives (FN).

The principal threshold-dependent quantities are:

\begin{description}
  \item[Sensitivity or recall.]
  The fraction of true ignitions correctly identified:
  \[
  \mathrm{sensitivity}
  =\mathrm{recall}
  =\frac{\mathrm{TP}}{\mathrm{TP}+\mathrm{FN}}.
  \]

  \item[Specificity.]
  The fraction of true non-ignitions correctly identified:
  \[
  \mathrm{specificity}
  =\frac{\mathrm{TN}}{\mathrm{TN}+\mathrm{FP}}.
  \]

  \item[False-positive rate.]
  The fraction of true non-ignitions incorrectly classified as ignition:
  \[
  \mathrm{FPR}
  =1-\mathrm{specificity}
  =\frac{\mathrm{FP}}{\mathrm{FP}+\mathrm{TN}}.
  \]

  \item[Precision.]
  The fraction of predicted ignitions that actually ignite:
  \[
  \mathrm{precision}
  =\frac{\mathrm{TP}}{\mathrm{TP}+\mathrm{FP}}.
  \]
\end{description}

Because the assembled ignition dataset is majority-positive, accuracy is not a
defensible selection metric. A model that predicts ignition for every
observation would already achieve a high accuracy while providing no useful
discrimination between ignition and no-ignition conditions.

Model comparison therefore begins with threshold-independent metrics:

\begin{description}
  \item[ROC--AUC.]
  The receiver-operating-characteristic curve plots sensitivity against the
  false-positive rate across all possible thresholds. ROC--AUC is the
  probability that a randomly selected ignition receives a higher predicted
  probability than a randomly selected non-ignition. A no-skill ranking model
  has ROC--AUC equal to \(0.5\).

  \item[PR--AUC.]
  The precision--recall curve plots precision against recall across all
  thresholds. Its no-skill baseline is the positive-class prevalence rather
  than \(0.5\). Since ignition is the majority label in this database,
  PR--AUC measures whether the model provides useful positive-class ranking
  beyond an always-ignite policy.

  \item[Brier score.]
  The Brier score measures squared probability error:
  \[
  \mathrm{BS}
  =\frac{1}{n}\sum_{i=1}^{n}
  \left(\widehat p_i-y_i\right)^2.
  \]
  Lower values indicate predicted probabilities that are closer to the
  observed labels.
\end{description}

Four threshold-selection objectives are also evaluated:

\begin{description}
  \item[F1 score.]
  The harmonic mean of precision and recall:
  \[
  \mathrm{F1}
  =\frac{2\,\mathrm{precision}\cdot\mathrm{recall}}
  {\mathrm{precision}+\mathrm{recall}}.
  \]

  \item[Balanced accuracy.]
  The average of sensitivity and specificity:
  \[
  \mathrm{balanced\ accuracy}
  =\frac{\mathrm{sensitivity}+\mathrm{specificity}}{2}.
  \]

  \item[Youden's \(J\).]
  The difference between sensitivity and false-positive rate:
  \[
  J=\mathrm{sensitivity}+\mathrm{specificity}-1.
  \]

  \item[Matthews correlation coefficient (MCC).]
  A balanced summary based on all four confusion-matrix counts:
  \[
  \mathrm{MCC}
  =\frac{\mathrm{TP}\,\mathrm{TN}-\mathrm{FP}\,\mathrm{FN}}
  {\sqrt{(\mathrm{TP}+\mathrm{FP})(\mathrm{TP}+\mathrm{FN})
  (\mathrm{TN}+\mathrm{FP})(\mathrm{TN}+\mathrm{FN})}}.
  \]
\end{description}

ROC--AUC, PR--AUC, and Brier score are used for candidate comparison because
they do not depend on a single operating threshold. Within each outer training
partition, four thresholds are selected from inner out-of-fold predictions:
the thresholds that maximize MCC, F1, balanced accuracy, and Youden's \(J\).
These thresholds are then frozen and applied once to the untouched outer test
partition. Outer-test labels are never used to choose a threshold.

Thresholded metrics are reported to describe possible operating points, not to
claim a single universally optimal fire-safety threshold. The preferred
operating point depends on the relative consequences of missing an ignition and
raising a false alarm.

\subsection{Candidate design}

The ignition portfolio contains 26 candidates across five families
(\cref{fig:method-ignition-xgb} and  \cref{fig:method-ignition-portfolio}). The design is an ablation of
how a family is trained, not a search over dozens of unrelated
architectures. Every family includes an unweighted baseline. XGBoost, as the
primary model, receives the full imbalance grid; the other families act as benchmarks.

Two imbalance mechanisms are treated as orthogonal. Paper weighting corrects
sampling bias from large campaigns. Class weighting rescales the loss by label
frequency. Focal loss \cite{lin2017focal} down-weights easy, confident
examples dynamically and is therefore an alternative class-imbalance
mechanism, not a third knob to stack on class weights. Focal candidates use
\(\gamma=2\) and \texttt{class\_weight:false} to avoid double-correcting the
majority class. Paper weighting is retained with focal loss because it
corrects a different bias.

\begin{figure}[htbp]
\centering
\includegraphics[width=0.9\textwidth]{ignition_xgb.png}
\caption{The primary XGBoost
family tests paper weighting, class weighting, and focal loss as separate
responses to paper-size and class imbalance. Total is 8 candidates}
\label{fig:method-ignition-xgb}
\end{figure}

The remaining 18 candidates test one structural hypothesis per family
(\cref{fig:method-ignition-portfolio}). Decision-tree depth, leaf size, and
split criterion are already inner-search hyperparameters, so no extra
``pruned tree'' candidate is added. Distance-weighted \(k\)-NN lets closer
experimental neighbours vote more strongly. Linear SVM tests whether the
ignition boundary is linearly separable in the physics feature space, which is
plausible if oxygen and pressure thresholds dominate. MLP focal variants
mirror the XGBoost focal logic at reduced scale.

\begin{figure}[htbp]
\centering
\includegraphics[width=0.86\textwidth]{ignition_portfolio.png}
\caption{Other ignition candidate portfolio. Decision tree, \(k\)-NN, MLP, and SVM
candidates provide structured comparison families. Total is 18 candidates.}
\label{fig:method-ignition-portfolio}
\end{figure}

Hyperparameters that are not part of a candidate identity are sampled
inside each outer training partition (\cref{tab:ignition-search}). Inner
ranking uses mean ROC--AUC, then PR--AUC, then lower Brier score.

\begin{table}[htbp]
\centering
\small
\caption{Inner-search spaces for ignition families}
\label{tab:ignition-search}
\begin{tabular}{@{}lp{0.72\textwidth}@{}}
\toprule
Family & Sampled hyperparameters \\
\midrule
XGBoost & \(n_{\mathrm{estimators}}\in\{200,400,600,800\}\); learning rate
\(\{0.01,0.03,0.05,0.1\}\); depth \(\{2,\ldots,6\}\); min-child-weight,
subsample, column sample, \(\ell_1/\ell_2\) regularisation, gamma. \\
Decision tree & depth \(\{3,5,7,10,\infty\}\); min-leaf
\(\{1,2,5,10,20\}\); split criterion gini, entropy, or log-loss. \\
\(k\)-NN & \(k\in\{5,9,15,25,40,60\}\); uniform or distance weights unless
fixed; Minkowski \(p\in\{1,2\}\); leaf size. \\
MLP & hidden widths \((32)\), \((64)\), \((128,64)\), \((128,64,32)\); \(\ell_2\)
and learning-rate grid; ReLU; batch size 32 or 128. \\
SVM & \(C\in\{0.1,0.3,1,3,10,30\}\); RBF \(\gamma\) grid or linear kernel when
fixed; shrinking on/off. \\
\bottomrule
\end{tabular}
\end{table}
\subsection{Pipeline procedure}

The ignition pipeline follows the same ordered procedure for every candidate:

\begin{enumerate}
  \item The cleaned ignition dataset and fixed feature manifest are loaded.
  Paper identity is retained only for grouping, weighting, and integrity
  checks; it is not used as a predictor.

  \item Outer evaluation assignments are frozen before candidate comparison.
  Four protocols are stored:
  \begin{itemize}
    \item repeated five-fold stratified interpolation cross-validation;
    \item a stratified 80/20 holdout used only as a consistency check;
    \item repeated five-fold paper-grouped extrapolation;
    \item leave-one-paper-out (LOPO) evaluation used as a robustness analysis
    for the extrapolation champion.
  \end{itemize}

  \item For one candidate and one outer split, only the outer training
  partition is used to fit preprocessing and evaluate hyperparameter
  configurations through inner cross-validation.

  \item Inner out-of-fold predicted probabilities are used to select four
  threshold policies: maximum MCC, maximum F1, maximum balanced accuracy, and
  maximum Youden's \(J\). The selected thresholds are frozen before the outer
  test partition is accessed.

  \item The selected hyperparameter configuration is refitted using the full
  outer training partition. It then produces probabilities for the untouched
  outer test partition.

  \item Threshold-independent metrics, including ROC--AUC, PR--AUC, and Brier
  score, are calculated from the outer-test probabilities. Threshold-dependent
  metrics are calculated using the thresholds selected within training.

  \item Outer-fold predictions, metrics, selected hyperparameters, threshold
  values, and integrity checks are persisted.

  \item After all eligible candidates are evaluated, the selection policy
  chooses separate interpolation and extrapolation champions. Only those
  selected configurations are refitted on all labelled ignition rows for
  inference.
\end{enumerate}

\subsection{Interpolation protocol}

Ignition interpolation uses five-fold stratified cross-validation repeated for
seeds 42--44, giving 15 outer folds. Stratification keeps the ignition/no-
ignition proportion similar in each fold while rows are assigned independently,
so papers may appear on both sides. Within each outer training partition,
stratified inner folds search the family space. A separate stratified 80/20
holdout for each seed is a quick consistency check, not the primary selection
protocol.

Interpolation selection requires valid probabilities, complete folds, passed
integrity checks, and ROC--AUC fold standard deviation no greater than 0.15.
Candidates within 0.01 of the best ROC--AUC are resolved by higher PR--AUC,
lower uncertainty, and a declared simplicity order (decision tree, \(k\)-NN,
SVM, XGBoost, MLP). The resulting champion is refitted on all
ignition-labeled rows for within-domain inference.

\subsection{Extrapolation and leave-one-paper-out}

Ignition extrapolation uses \texttt{StratifiedGroupKFold}: all rows from one
paper stay together while each fold retains both classes where feasible. Five
paper-grouped folds repeated across three seeds again produce 15 outer tests.
The inner search is also paper-disjoint, so a paper held out for inner
validation cannot appear in the corresponding inner training fold.

The primary metric remains ROC--AUC, but eligibility additionally requires a
successful LOPO analysis. Each of the 87 papers is held out once. A
single-class paper cannot define its own ROC--AUC, so LOPO is interpreted
through pooled predictions, calibration, and the subset of papers for which a
metric is defined. The selected extrapolation model and its thresholds are
refitted on the full labeled cohort only after this paper-disjoint comparison
is frozen.

\section{Flame-spread-rate regression pipeline}
\label{sec:regression-method}
\subsection{Task and metrics}

The FSR pipeline predicts a continuous flame spread rate in
\si{\metre\per\second}. Unlike ignition classification, regression does not
require a probability threshold: each prediction is compared directly with the
measured FSR value.

For observation \(i\), let \(y_i\) denote the measured FSR and
\(\widehat y_i\) the predicted FSR. The residual is

\[
e_i=\widehat y_i-y_i.
\]

A positive residual indicates overprediction, whereas a negative residual
indicates underprediction. The pipeline reports several complementary metrics:

\begin{description}
  \item[Root-mean-square error (RMSE).]
  \[
  \mathrm{RMSE}
  =\sqrt{\frac{1}{n}\sum_{i=1}^{n}
  \left(y_i-\widehat y_i\right)^2}.
  \]
  RMSE is expressed in \si{\metre\per\second} and penalizes large prediction
  errors more strongly than small errors. It is the primary model-selection
  metric because it directly represents the physical size of the prediction
  error.

  \item[Mean absolute error (MAE).]
  \[
  \mathrm{MAE}
  =\frac{1}{n}\sum_{i=1}^{n}
  \left|y_i-\widehat y_i\right|.
  \]
  MAE is also expressed in \si{\metre\per\second}. It represents the typical
  absolute prediction error and is less sensitive than RMSE to a small number
  of large misses.

  \item[Coefficient of determination (\(R^2\)).]
  \[
  R^2
  =1-\frac{\sum_{i=1}^{n}(y_i-\widehat y_i)^2}
  {\sum_{i=1}^{n}(y_i-\bar y)^2}.
  \]
  \(R^2\) compares the model with a predictor that always returns the mean
  target value. A value of \(R^2=1\) is a perfect fit, whereas \(R^2=0\)
  performs no better than the mean predictor. Because \(R^2\) depends on the
  target variance in the evaluated fold, it can differ substantially across
  held-out paper groups even when absolute errors are similar.

\end{description}

All predictions are transformed back to the original FSR scale before scoring.
This includes candidates trained with the signed-log target transformation.
Consequently, every reported RMSE and MAE values remains interpretable in
physical FSR units.

RMSE and \(R^2\) answer complementary questions. RMSE is the preferred metric
for an engineering error budget because it is reported in
\si{\metre\per\second}. \(R^2\) indicates whether the model explains
variation beyond a mean predictor, but it is sensitive to the target spread of
the held-out observations. The two metrics should therefore be interpreted
together rather than treated as interchangeable.

Interpolation and paper-grouped extrapolation also answer different questions.
Interpolation measures performance on new rows from papers that may otherwise
appear in training. Paper-grouped extrapolation measures transfer to entirely
unseen publications or campaigns. These scores must not be averaged or
substituted for one another; the paper-grouped result is the primary evidence
for transfer beyond represented source papers.

Jose's headline \(R^2>0.9\) values were obtained on a narrower wire dataset
under random ten-fold cross-validation \cite{riveraWireML}. The present
pipeline reports an analogous interpolation result for comparison, but its
scientific generalization claim is based on the paper-grouped evaluation.

\subsection{Candidate design}

The FSR portfolio contains 34 candidates across five families
(\cref{fig:method-fsr-xgb} and \cref{fig:method-fsr-portfolio} ). Each candidate tests one hypothesis:
does paper weighting improve transfer; does a robust loss help with
experimental outliers; does compressing a skewed target help optimization;
is local similarity enough; is the physics feature space already linear?
Jose's linear families are omitted. Random forest is added as the bagging
counterpart to XGBoost so that any later preference for trees is not
confounded with boosting.

XGBoost (10 candidates) is again the most granular family. Three objectives
are compared under square-root and inverse paper weights: squared error,
absolute error, and pseudo-Huber. Absolute error reduces sensitivity to large
residuals. Pseudo-Huber is smooth near zero and linear in the tails. Two
further candidates predict \(\mathrm{sign}(y)\log(1+|y|)\) and invert the
transform at inference, testing whether compressing the heavy FSR tail
improves optimization. One paper-bagging candidate averages ten XGBoost
models trained on paper-level bootstrap samples, reducing dependence on any
single campaign.

Random forest (7) mirrors the scientifically interesting XGBoost knobs without
boosting: baseline, sqrt/inverse paper weights, absolute-error split
criterion, and signed-log targets. \(k\)-NN (6) crosses uniform versus
distance neighbours with none/sqrt/inverse paper weights. MLP (5) includes
baseline and paper-weighted models plus two signed-log networks, because
neural optimization is often more stable on compressed targets when the
response spans orders of magnitude. SVR (6) crosses linear versus RBF kernels
with the same three paper weights: a strong linear kernel would imply that
feature engineering already captured most of the spread physics.


\begin{figure}[htbp]
\centering
\includegraphics[width=0.78\textwidth]{fsr_xgb.png}
\caption{XGBoost FSR candidate-design ablations. The XGBoost family tests
paper weighting, robust objectives, signed-log targets, and paper-level
bagging as distinct modelling hypotheses. Total is 10 candidates.}
\label{fig:method-fsr-xgb}
\end{figure}


\begin{figure}[htbp]
\centering
\includegraphics[width=0.86\textwidth]{fsr_portfolio.png}
\caption{Complete FSR candidate portfolio across Random Forest,
\(k\)-NN, MLP, and SVR. Candidate variants test weighting, robustness, target
representation, local similarity, bagging, and kernel assumptions. Total is 24 candidates.}
\label{fig:method-fsr-portfolio}
\end{figure}

Inner search again samples family hyperparameters rather than candidate
identity (\cref{tab:fsr-search}). Configurations are ranked by mean inner RMSE,
then MAE, then higher \(R^2\).

\begin{table}[htbp]
\centering
\small
\caption{Inner-search spaces for FSR families.}
\label{tab:fsr-search}
\begin{tabular}{@{}lp{0.72\textwidth}@{}}
\toprule
Family & Sampled hyperparameters \\
\midrule
XGBoost & \(n_{\mathrm{estimators}}\in\{100,200,300,500,800\}\); learning rate
\(\{0.01,0.02,0.05,0.1,0.2\}\); depth \(\{2,\ldots,6\}\); min-child-weight,
subsample, column sample, \(\ell_1/\ell_2\), gamma. \\
Random forest & 100--500 trees; depth 4--12 or unlimited; split
and leaf minima; \texttt{max\_features} in \(\{\sqrt{\,},\,\log_2,0.5,0.8,1\}\). \\
\(k\)-NN & \(k\in\{3,5,7,9,11,15,21,31\}\); uniform or distance unless fixed;
\(p\in\{1,2\}\). \\
MLP & widths including \((100)\), \((128,64)\), \((256,128)\); ReLU or tanh;
\(\ell_2\) and learning-rate grid; batch size 32--128. \\
SVR & \(C\in\{0.1,\ldots,100\}\); \(\varepsilon\in\{0.01,\ldots,0.5\}\); RBF
\(\gamma\) grid or linear kernel when fixed. \\
\bottomrule
\end{tabular}
\end{table}

\subsection{Pipeline procedure}

The FSR pipeline evaluates every candidate under two training-data conditions:
raw data only (\texttt{rd}) and raw data plus bootstrap augmentation
(\texttt{rd\_bt}). The procedure is:

\begin{enumerate}
  \item The cleaned FSR dataset is loaded. Only rows with a numeric FSR target
  are retained. The target is measured and reported in
  \si{\metre\per\second}.

  \item Outer evaluation assignments are frozen before candidate comparison.
  Interpolation uses repeated approximately 80/20 row-level holdouts. Where
  feasible, FSR values are binned into six target quantiles during split
  construction so that the strongly right-skewed target distribution is
  represented in both partitions.

  \item Extrapolation also uses repeated approximately 80/20 holdouts, but the
  assignment unit is the complete canonical paper. For each repeat, candidate
  paper-disjoint partitions are generated and the most target-balanced
  partition is retained. No paper can appear in both the corresponding
  training and test partitions.

  \item For each candidate and outer split, hyperparameters are selected using
  only the outer training partition. Interpolation uses shuffled inner folds;
  extrapolation uses paper-disjoint inner folds.

  \item The \texttt{rd} condition fits the selected configuration using the
  raw outer-training rows. The \texttt{rd\_bt} condition appends 1,000 rows
  sampled with replacement from those same training rows. Each copied row keeps
  its original predictors, FSR target, and source-paper identity. Test rows are
  never augmented.

  \item Each fitted model predicts the untouched outer test partition. RMSE,
  MAE, \(R^2\), mean bias error, and other diagnostic metrics are calculated
  after predictions are converted back to the original FSR scale.

  \item The pipeline persists outer-test predictions, fold metrics, selected
  hyperparameters, inner-search histories, and integrity records for every
  eligible candidate, protocol, and augmentation condition.

  \item Candidate shards are merged after evaluation. Selection first
  prioritises low mean RMSE, then MAE, higher \(R^2\), uncertainty, preference
  for raw data when candidates are tied, and the declared model-simplicity
  order. The selected interpolation and extrapolation configurations are then
  refitted on all labelled FSR rows for inference.
\end{enumerate}

\subsection{Interpolation protocol}

FSR interpolation uses ten repeated, approximately 80/20 row holdouts with
seeds 42--51. Target values are binned into six quantiles where feasible so
that each split represents the strongly skewed FSR distribution. Papers may
occur on both sides. Inside each outer training partition, shuffled \(K\)-fold
validation compares hyperparameters. The selected recipe is fitted on all
2,024 training rows and tested on the 507 untouched rows for that repeat.

The selection policy rejects incomplete or non-finite evaluations, RMSE
coefficient of variation above 0.45, and mean \(R^2<-0.25\). Candidates within
2\% of the lowest RMSE form a tie band resolved by higher \(R^2\), lower
uncertainty, preference for real data over bootstrap copies, simpler family
(random forest, \(k\)-NN, SVR, XGBoost, MLP), and fewer features where the
feature sets differ.

\subsection{Extrapolation protocol}

FSR extrapolation also uses ten repeated 80/20 holdouts, but the unit assigned
to train or test is the paper. For each repeat, 50 candidate
\texttt{GroupShuffleSplit} partitions are generated and the most
target-balanced paper-disjoint split is retained. Each test contains 14 papers
and each training partition 52; row counts vary with paper size. Balance
screening reduces avoidable target-range changes without relaxing paper
separation. Because the screening criterion uses only target-distribution
summaries to choose among pre-generated paper partitions, it is treated as part
of split construction and is frozen before any candidate is evaluated. The
result remains a transfer estimate for this balance-screened split policy,
not for an arbitrary future collection of papers.

Inner tuning uses \texttt{GroupKFold} by paper. Extrapolation tie-breaks add
the interpolation--extrapolation gap so that a recipe is not favoured solely
for excellent within-paper fit. Cluster bootstrap intervals resample papers
rather than thousands of correlated rows. The winning
\((\mathrm{candidate},\,\mathrm{augmentation})\) pair is refitted on all
2,531 FSR rows and stored separately from the outer-fold evidence.